\section{Introduction}
\label{sec:v75-introduction}
\label{sec:v76-introduction}
\label{sec:v70-introduction}

A long billiard trajectory constrained to follow a hyperbolic periodic
orbit has exponentially small phase volume.  Its endpoints nevertheless
range over fixed neighborhoods of the reflecting contacts.  We ask
whether the conditional distribution of these endpoints determines the
reflecting arcs as smooth functions.  Recovering their Taylor series is
not sufficient: a nonzero smooth perturbation can be flat at every
marked contact.

We prove that two positive excess-time laws at each phase determine the
complete smooth contact germs.  The polygon and tangent frames are
supplied, but the curvatures, graph heights and local flux amplitudes are
not.  The argument has two substantive steps.  First, an exact Jacobi
normalization and a two-ended determinant comparison preserve the
nonlinear physical law at its exponentially small scale.  Second, the
stationary envelope on actual reflecting graphs separates the boundary
value at the initial contact from all later, contracted visits.  A global
positive quadratic inverse and a finite signed jet alignment make this
second step applicable without analyticity or a closeness assumption.

\subsection{The experiment and the main theorem}

Fix a clear nongrazing periodic polygon with $r\ge2$ distinct physical
contacts in a planar periodic dispersing billiard.  Its lifted vertices
$q_i$, signed unit tangents $t_i$, outward obstacle normals $n_i$, phase
labels and closing translation are fixed marks.  The reference polygon
satisfies specular reflection.  Write
\[
 L_i=|q_{i+1}-q_i|,\quad e_i=(q_{i+1}-q_i)/L_i,\quad
 \nu_i=e_i\cdot n_i>0,\quad p_i=e_i\cdot t_i/\nu_i.
\]
The unknown arcs are the anchored normal graphs
\begin{equation}\label{eq:v75-graphs}
 R_i(x)=q_i+\nu_i^{-1}\{t_i x-n_iF_i(x)\},\qquad
 F_i(0)=F_i'(0)=0,\qquad F_i''(0)>0.
\end{equation}
The observed coordinate $x=\nu_i t_i\cdot(R_i-q_i)$ is a tangent
projection.  It reveals neither $F_i(x)$ nor an arclength conversion.
Clearance means that the open reference flights meet no unintended
obstacle; the precise uniform margins are given in
Section~\ref{sec:v64-periodic-itinerary}.

Let $P$ be a returning oriented period, doubled when necessary to return
the signed frames, and let $\mathcal L$ be its length.  For $N=nP$
flights starting at phase $b$, let $W_{N,b}(u,v)$ be the stationary
bridge length in the prescribed collars.  Set
\[
 E_{N,b}=W_{N,b}-n\mathcal L+p_bu-p_bv,\qquad
 D_{N,b}=-W_{N,b,uv}(0,0)>0.
\]
At time $n\mathcal L+d$, success requires the specified first collision,
these $N$ flights in the itinerary collars, no further collision, and
endpoints in one fixed small square $J^2$.  Here $d>0$ is fixed as
$N$ increases.  The square is chosen so that the residual time is
positive and smaller than either adjacent free-flight length.

Preparations are independent draws from normalized unit-speed Liouville
measure on the full phase space.  Itinerary selection, gate membership
and success/failure tags are exact.  The conditional-law theorem does
not require the free-area normalizer; its reconstruction from counts
below does require this specified preparation measure.  Bounded sensor
errors in the finite experiment affect accepted tangent-coordinate
readouts only.

\begin{theorem}[Relative boundary laws and smooth contact determination]
\label{thm:v75-main}
\label{thm:v76-main}\label{thm:v70-main}
For a marked polygon as above with positive-curvature $C^\infty$
contacts, the following assertions hold on a contact collar independent
of $N$.

There are anchored smooth half-line actions $S_b^-,S_b^+$ and positive
smooth amplitudes $B_b^-,B_b^+$, normalized by $B_b^\pm(0)=1$, such
that, for each separately fixed finite order $k$,
\begin{equation}\label{eq:v75-relative}
 \|E_{N,b}-S_b^-\oplus S_b^+\|_{C^k}
 +\|(-W_{N,b,uv})/D_{N,b}-B_b^-\otimes B_b^+\|_{C^k}
 \le C_k\tau^N,\qquad 0<\tau<1.
\end{equation}
The reference is exact:
\begin{equation}\label{eq:v75-reference}
 D_{nP,b}=\frac{\sinh\chi}
  {(\mathcal M_b)_{12}\sinh(n\chi)},\qquad
 \operatorname{tr}\mathcal M_b=2\cosh\chi.
\end{equation}
The monodromy $\mathcal M_b$ is computed from the actual geometric
Jacobi operator; it is not additional observation data.

For two fixed offsets $0<d_1<d_2$ in the positive physical window,
the conditional densities on $J^2$ converge in each fixed $C^k$ norm
to
\begin{equation}\label{eq:v75-law}
 f_{b,d}(u,v)=Z_{b,d}^{-1}B_b^-(u)B_b^+(v)
       \{d-A_b(u)-C_b(v)\},\qquad d\in\{d_1,d_2\},
\end{equation}
where $A_b=S_b^--p_b\,\mathrm{id}$ and
$C_b=S_b^++p_b\,\mathrm{id}$.  The positive constants $Z_{b,d}$
are not supplied.  Equality of these phase-resolved laws for two actual
contact tuples implies equality of their smooth contact germs in
\eqref{eq:v75-graphs}.  No curvature input, reflection symmetry,
analyticity or closeness between the candidates is required.

On a fixed bounded positive class, an integer $m$ and a common smaller
collar $I_c$ can be chosen so that the inverse also satisfies
\begin{equation}\label{eq:v75-stability}
 \|F-G\|_{C^0(I_c)}
 \le C\max_{b,d}\|f^F_{b,d}-f^G_{b,d}\|_{C^m(I^2)}
\end{equation}
for sufficiently small right side, with $I_c\Subset I\subset J$.
The class has the $C^{m+3}$ graph bounds and positive quadratic bounds
of \eqref{eq:v68-smooth-class}; the densities retain their normalization
on $J^2$ when restricted to $I^2$.  The order $m$ is chosen from the
contraction margin, not uniformly over degenerating geometric classes.
\end{theorem}
\begin{proof}
The relative and phase-volume statements are
Theorems~\ref{thm:v64-periodic-relative} and
\ref{thm:v64-physical-law}, transferred to measured graph coordinates
by Lemma~\ref{lem:v65-signed-jacobi}.
Proposition~\ref{prop:v65-two-offset} recovers the actions.
Theorem~\ref{thm:v66-curvature} identifies their positive quadratic
parameters globally, and Theorem~\ref{thm:v65-cyclic-jets} aligns
finite jets.  The actual-function estimate in
Proposition~\ref{prop:v68-weighted-separation} proves smooth uniqueness
and, after polynomial alignment, Theorem~\ref{thm:v68-smooth-stability}
proves \eqref{eq:v75-stability}.  All these proofs are included below.
\end{proof}

\subsection{The two analytic steps}

An absolute estimate for the stationary action does not control a
conditional law whose mass is of order $D_{N,b}$.  The exact
unconditioned density, with free cell area $\mathcal A$, is
\begin{equation}\label{eq:v75-physical}
 \frac{D_{N,b}}{2\pi\mathcal A}\,
 \frac{-W_{N,b,uv}(u,v)}{D_{N,b}}
       \{d-E_{N,b}(u,v)+p_bu-p_bv\}.
\end{equation}
Thus the required estimate is relative before conditioning, and the
linear endpoint momenta must be restored before the clock is read.
The cofactor formula cancels $D_{N,b}$ exactly.  Its logarithm contains
a sum of local mixed-chord terms and
$\log\det(I+G_N\Delta H_N)$, where $G_N=H_N(0)^{-1}$.
Although the number of interior sites increases, endpoint localization
keeps $\Delta H_N$ uniformly trace class.  Separating two end blocks,
controlling the opposite-end Green entries and differentiating the
convergent logarithmic series prove \eqref{eq:v75-relative}.
A quadratic approximation, or an absolute action error divided by the
success probability, would not provide this estimate.

The action extraction is elementary once this physical limit is known.
At one phase set
\[
 Q(u,v)=\frac{f_{d_1}(u,v)f_{d_2}(0,0)}
                 {f_{d_2}(u,v)f_{d_1}(0,0)}.
\]
Then
\[
 \frac{d_1d_2(1-Q)}{d_2-d_1Q}=A(u)+C(v).
\]
The denominator has a positive floor on the common gate.  Axis evaluation
recovers the two actions without knowing either flux amplitude.

The passage from actions to graphs is functional, not formal.
The quadratic Schur complement gives a global inverse on the positive
curvature image.  With these curvatures fixed, the signed cyclic
response at each higher order is invertible.  This identifies the finite
jets required by the next argument, but it cannot itself exclude a flat
boundary difference.  For two actual graph tuples, integrate the
stationary envelope along their local interpolation to obtain
\begin{equation}\label{eq:v75-envelope}
 \mathcal S(G)-\mathcal S(F)=\overline\alpha(G-F)+K(G-F),
 \qquad \overline\alpha\ge\tfrac12.
\end{equation}
Every summand of $K$ evaluates a boundary difference at a later visit
$x_{b,j}(u)$ satisfying $|x_{b,j}(u)|\le a^j|u|$ with $a<1$ and
multiplicative constant one.  For differences vanishing through order
$m-1$,
\[
 \|h\|_{m,R}=\max_b\sup_{0<|u|\le R}\frac{|h_b(u)|}{|u|^m},
 \qquad
 \|\overline\alpha^{-1}K\|\le\frac{6a^m}{1-a^m}<1
\]
after choosing a finite $m$.  This separates actual smooth functions,
including their flat parts.  The terminal term of the finite stationary
envelope is estimated before the half-line limit is taken.  Polynomial
jet alignment makes the same estimate quantitative for unequal data.

Proposition~\ref{prop:v68-flat-family} realizes the distinction on
closed obstacles: an $\varepsilon\chi(u)e^{-1/u^2}$ contact
perturbation, compensated by a remote area adjustment, preserves every
contact and action Taylor jet but changes the complete boundary laws.
The local response is not manufactured by a change of normalization.
The complementary locality result in Section~\ref{sec:v71-local-observation}
identifies what the observation leaves free away from the visited arcs.

The two estimates admit a formulation independent of the billiard
notation.  Proposition~\ref{prop:v76-determinant} isolates the relative
comparison of two localized end blocks: the first off-diagonal variation
of the logarithmic determinant vanishes, so their coupling enters
quadratically in Hilbert--Schmidt norm, with no dimension factor.
Proposition~\ref{prop:v76-envelope} bounds an actual integrated envelope
by a nonnegative visit matrix and inverts it when its spectral radius is
less than one.  These statements express the analytic mechanism used
here; they do not identify different observation models.

For periodic visits the latter estimate has a quantitative refinement.
If $\widehat a_i<1$ are common one-step bounds on the interpolation
class, set $\overline a=(\prod_i\widehat a_i)^{1/r}$.
Corollaries~\ref{cor:v76-phase-separation} and
\ref{cor:v76-phase-stability} replace the sufficient scalar condition
$6a^m/(1-a^m)<1$ by $6\overline a^m/(1-\overline a^m)<1$.
Their explicit phase weights carry a condition-number cost in the
unweighted estimate.  This improves the sufficient order criterion on
heterogeneous classes, not the target or the information supplied to
the inverse.  It need not improve every conditioning constant.

\subsection{What the observation refinements add}

The remaining local results use this same inverse rather than supplying
additional rigidity mechanisms.  At oblique phases, the marked nonzero
slope identifies an unreported positive offset from one law; the
fixed-slice extraction in Section~\ref{sec:v72-one-law} avoids an
additional derivative at each inverse order.  For normalized Liouville
preparations, the same-gate success count identifies the free-area
scalar.  Section~\ref{sec:v73-complete-record} proves both directions
of the resulting invariant: contact germs, offsets and free area.
The exact area identity uses the actual finite density, and its
stability retains the factor $1+N$ from logarithmic twist sensitivity.
It does not determine unvisited boundary arcs.

Section~\ref{sec:v74-moments} recovers the same limiting density from
four one-dimensional weighted profiles through $MH^{-1}N$.
The cancellation is the classical low-rank mechanism, as in
\cite{HammHuang2019}; the application requires signed geometric
nondegeneracy and the finite-flight defect.  The profiles retain paired
endpoint information.  They are not exact sufficient statistics for an
arbitrary finite-flight density.  On a fixed oblique gate they reduce
the smoothing dimension, giving the sufficient exponents
\[
 \alpha_1=\frac{s}{2(m+s)+1},\qquad
 \gamma_1=\frac{s}{m+s+2},\qquad
 \beta_1=\frac{\alpha_1\omega}{\omega+\alpha_1\Gamma}.
\]
With the explicit count, bandwidth and small-error conditions of
Theorem~\ref{thm:v74-charged} and Corollary~\ref{cor:v74-budget},
$rB$ charged preparations give profile-and-offset error $Ct$ and joint
logarithmic-area error $Ct\log(e/t)$ with probability at least
$1-2\zeta$, where $t=(L_B/B)^{\beta_1}+\delta^{\gamma_1}$.
These are upper bounds for a measurable actual-table estimator, not
claims of minimax optimality or effective enumeration.
The affine example in Remark~\ref{rem:v75-conditioning} exhibits the
loss of moment conditioning as the gate or obliquity shrinks.
Section~\ref{sec:v76-conditioning} combines this with the inverse order,
phase weights, offset separation, rare acceptance and finite-area
sensitivity.  These sufficient estimates are not efficiency guarantees.

\subsection{Relation to other inverse problems}
\label{sec:v71-lens-comparison}

The corrected planar theorem of Noakes--Stoyanov
\cite{NoakesStoyanov2020} globally determines finite unions of strictly
convex $C^3$ obstacles from exterior travelling-time or scattering-length
data.  Smooth obstacle determination from rich exterior data is therefore
not new in itself.  Here the table is periodic, an internal polygon and
its frames are supplied, and the observations are localized conditional
endpoint laws.  Our smooth target is the visited contact germs.
Equal-area remote completions in
Corollary~\ref{cor:v71-local-completion} have the same local record;
their exterior scattering data on a common compact realization differ.
This explains the spatial scope without identifying the two data sets.

De Simoi--Kaloshin--Leguil \cite{DSKL} use marked lengths for analytic
open billiards under symmetry and genericity assumptions.
Finamore--Leguil \cite{FinamoreLeguil} use an enriched marked length
spectrum for finite-horizon Sinai billiards.  Our result neither removes
hypotheses from those problems nor establishes same-data dominance.
Its content is the passage from a relatively controlled rare physical
law to an inverse sensitive to actual flat boundary changes.
The matrix identity, quotient extraction, concentration and
minimum-distance selection are supporting tools; they are not offered
as new general principles of rigidity or statistics.

\subsection{Organization}

Sections~\ref{sec:v64-periodic-itinerary}--\ref{sec:v68-smooth-contact}
contain the complete central proof: the physical relative law, measured
coordinates and stationary envelope, global curvature identification,
and smooth uniqueness with stability.  Section~\ref{sec:v71-local-observation}
proves exact locality and constructs actual remote completions.
Sections~\ref{sec:v72-one-law}--\ref{sec:v74-moments} develop the
oblique single-law protocol, its area-augmented invariant, and the
moment reconstruction with finite-rank defect and charged
one-dimensional estimation.  Section~\ref{sec:v76-conditioning} gives
the integrated conditioning and budget analysis.
Appendix~\ref{sec:v69-sampled-contact}
provides the two-offset finite-preparation argument; it is also needed
when a marked phase has normal incidence.

The accompanying \emph{Full technical manuscript} contains the complete
analytic, finite-stopping, alternating-channel, intrinsic multichannel,
lattice, calibration and coarsening arguments.  References prefixed
$\mathrm F$ point to its numbered statements.  In the principal article
these references identify extensions or special-case comparisons, not
omitted premises of the central smooth theorem.  In particular, neither
analytic continuation nor a global matching hypothesis is used in that
proof.  The proof organization follows this dependency distinction rather
than the order in which observation refinements were obtained.
